\documentclass[11pt]{article}
\usepackage{amsmath, amsthm, amssymb}
\usepackage[margin=1.1in]{geometry}
\usepackage{hyperref}

\newtheorem{theorem}{Theorem}
\newtheorem{proposition}[theorem]{Proposition}
\newtheorem{corollary}[theorem]{Corollary}
\newtheorem{lemma}[theorem]{Lemma}
\theoremstyle{plain}
\newtheorem*{theorem*}{Theorem}
\theoremstyle{definition}
\newtheorem*{conjecture}{Conjecture}
\theoremstyle{remark}
\newtheorem*{remark}{Remark}

\newcommand{\Frob}{\operatorname{Frob}}
\newcommand{\Ind}{\operatorname{Ind}}
\newcommand{\Res}{\operatorname{Res}}
\newcommand{\Hom}{\operatorname{Hom}}
\newcommand{\ch}{\operatorname{ch}}
\newcommand{\triv}{\mathrm{triv}}
\newcommand{\fdeg}{\widetilde{f}}
\newcommand{\CC}{\mathbb{C}}
\newcommand{\ZZ}{\mathbb{Z}}

\title{$[2r-1]_t$-divisibility of $S_r$-isotypic multiplicities of\\$R_{(k^r)}$ for arbitrary $k$: the Molien bound suffices}
\author{Clio}
\date{2026-08-04}

\begin{document}
\maketitle

\begin{abstract}
Fix integers $r \ge 2$, $k \ge 2$ and set $n = rk$, $d = 2r-1$,
$\alpha = (k^r)$. Let $R_\alpha = R_n^{S_\alpha}$ be the parabolic
coinvariant algebra of $S_n$ (a graded $S_r$-module via the residual
block-permutation action of $H/S_\alpha \cong S_r$), and let
$m_\pi(t) = \sum_\lambda \fdeg_\lambda(t) \langle s_\pi[h_k], s_\lambda\rangle$
be the graded $S_r$-isotypic multiplicities (Theorem~D). We prove:

\begin{theorem*}
Assume $d = 2r-1$ is prime and $k \bmod d \in \{2, 3, \ldots, d-1\}$. Then
\[
[d]_t \mid m_\pi(t) \qquad \text{for every } \pi \vdash r.
\]
\end{theorem*}

This subsumes and strengthens the ``first-period'' theorem of
\cite{Clio-first-period} (which restricted to $2 \le k \le d-1$).
The proof is a single-line refinement of the Molien-Springer argument
of \cite{Clio-first-period}: the previously-used inequality
$m_d(\tilde\sigma h) = 0$ (uniform term-by-term vanishing) is replaced
by the weaker $m_d(\tilde\sigma h) < \lfloor n/d\rfloor$, which suffices
for the trace $\text{tr}(\tilde\sigma h \mid R_n; t)$ to vanish at $\zeta_d$.
The weaker inequality holds throughout the full ``allowed'' range
$k \bmod d \in \{2, \ldots, d-1\}$, so \emph{term-by-term} vanishing
extends to arbitrarily large $k$; no cancellation across
$h$ is needed.

The theorem covers every empirically verified positive case with prime
$d = 2r-1$ (through $r = 6$; see
\texttt{\~{}/projects/probes/2026-07-31-2r-1-divisibility/}).
It is verified computationally for the smallest new cases
$(r, k) \in \{(3,7), (3,8), (3,9), (3,12), (4,9)\}$
in \S\ref{sec:verify}.
\end{abstract}

\section{Setup and reduction to characters}
\label{sec:setup}

We recall the notation and reductions of \cite{Clio-first-period}.

Let $S_n$ act on the polynomial ring $\CC[x_1,\ldots,x_n]$ by permutation
and let $R_n = \CC[x_1,\ldots,x_n] / (\CC[x_1,\ldots,x_n]_+^{S_n})$
be the coinvariant algebra. For $\alpha = (k^r)$, $n = rk$, let
$S_\alpha = S_k^r \subseteq S_n$, $H = S_k \wr S_r$ its wreath normaliser,
and $R_\alpha = R_n^{S_\alpha}$ the parabolic coinvariant algebra. The
residual quotient $H / S_\alpha \cong S_r$ acts on $R_\alpha$ by graded
ring automorphisms (\cite[Thm 1(i)]{Clio-Wreath-general-r}).

For each $\pi \vdash r$, the graded $\pi$-isotypic multiplicity is
$m_\pi(t) = \sum_i t^i \dim \Hom_{S_r}(V_\pi, (R_\alpha)_i)$;
Theorem~D of \cite{Clio-Wreath-general-r} gives the closed form
$m_\pi(t) = \sum_\lambda \fdeg_\lambda(t) \langle s_\pi[h_k], s_\lambda\rangle$.
For $\sigma \in S_r$, write
\[
P_\sigma(t) := \sum_{\pi \vdash r} \chi^\pi(\sigma) \, m_\pi(t)
\]
for the graded character of $\sigma$ on $R_\alpha$; by orthogonality,
\[
m_\pi(t) = \frac{1}{r!} \sum_{\sigma \in S_r} \chi^\pi(\sigma) \, P_\sigma(t).
\]

Since $d$ is prime, $[d]_t = \Phi_d(t)$ is irreducible over $\ZZ[t]$, so
$[d]_t \mid m_\pi(t)$ iff $m_\pi(\zeta_d) = 0$, where $\zeta_d = e^{2\pi i/d}$.

\begin{lemma}[{\cite[Lem 4.1]{Clio-first-period}}]
\label{lem:reduce}
$m_\pi(\zeta_d) = 0$ for every $\pi \vdash r$ if and only if
$P_\sigma(\zeta_d) = 0$ for every $\sigma \in S_r$.
\end{lemma}

The theorem thus reduces to the following.

\begin{proposition}
\label{prop:main}
For $r \ge 2$, $d = 2r - 1$ prime, and $k \bmod d \in \{2, \ldots, d-1\}$:
$P_\sigma(\zeta_d) = 0$ for every $\sigma \in S_r$.
\end{proposition}

\section{The Molien bound}

For $\sigma \in S_r$ we write $\tilde\sigma \in S_n$ for its
block-permutation lift, viewing
$S_r \hookrightarrow N_{S_n}(S_\alpha)/S_\alpha$.

\begin{lemma}[{Molien formula for $P_\sigma$; \cite[Lem 4.3]{Clio-first-period}}]
\label{lem:molien}
For every $\sigma \in S_r$,
\[
P_\sigma(t) \;=\;
\frac{1}{|S_\alpha|} \sum_{h \in S_\alpha}
\frac{\prod_{i=1}^{n}(1 - t^i)}{\prod_j (1 - t^{c_j(\tilde\sigma h)})},
\]
where the inner product ranges over cycles of $\tilde\sigma h$ as an
element of $S_n$.
\end{lemma}

\begin{lemma}[{Vanishing order of individual terms; \cite[Lem 4.4]{Clio-first-period}}]
\label{lem:order}
For any $g \in S_n$, let $m_d(g) = \#\{j : d \mid c_j(g)\}$. Then
$\text{tr}(g \mid R_n; t) = \prod_{i=1}^n(1-t^i)/\prod_j(1 - t^{c_j(g)})$
has a zero of order exactly $\lfloor n/d\rfloor - m_d(g)$ at $t = \zeta_d$.
In particular, $\text{tr}(g \mid R_n; \zeta_d) = 0$ iff
$m_d(g) < \lfloor n/d\rfloor$.
\end{lemma}

\begin{lemma}[{Wreath cycle formula; \cite[Lem 4.5]{Clio-first-period},
\cite[Ch.~I App.~B]{Macdonald}}]
\label{lem:wreath}
Let $\sigma \in S_r$, $h = (h_1, \ldots, h_r) \in S_\alpha = S_k^r$, and
$\tilde\sigma h \in S_n$. For each cycle $(b_1\ b_2\ \cdots\ b_\ell)$ of
$\sigma$, form the companion permutation
$g_c := h_{b_\ell} h_{b_{\ell-1}} \cdots h_{b_1} \in S_k$. Then the cycles
of $\tilde\sigma h$ on the $\ell k$ elements of blocks $b_1, \ldots, b_\ell$
have lengths precisely $\{\ell m : m \text{ is a cycle length of } g_c\}$.
\end{lemma}

We now prove the main bound.

\begin{proposition}[Main bound]
\label{prop:bound}
Let $r \ge 2$, $d = 2r-1$ prime, $k \ge 2$, and $\sigma \in S_r$ with
$p = p(\sigma)$ cycles (counting fixed points). Then
\[
\max_{h \in S_\alpha} m_d(\tilde\sigma h) \;\le\; p \cdot \lfloor k/d\rfloor.
\]
\end{proposition}

\begin{proof}
By Lemma \ref{lem:wreath}, every cycle of $\tilde\sigma h$ has length
$\ell m$ for some cycle length $\ell$ of $\sigma$ (with $\ell \le r$) and
some cycle length $m$ of the corresponding companion $g_c \in S_k$.
Since $d = 2r - 1 > r \ge \ell$ and $d$ is prime, $\gcd(\ell, d) = 1$;
hence $d \mid \ell m$ if and only if $d \mid m$. Consequently
\[
m_d(\tilde\sigma h) \;=\; \sum_{\text{cycles $C$ of }\sigma} m_d(g_c(C)),
\]
where $g_c(C) \in S_k$ is the companion of cycle $C$.

For any $g \in S_k$, cycles of $g$ of length divisible by $d$ each use at
least $d$ elements out of $k$, so $m_d(g) \le \lfloor k/d\rfloor$. Summing
over the $p$ cycles of $\sigma$ gives the claim.
\end{proof}

\begin{corollary}
\label{cor:gap}
Write $k = ad + b$ with $a = \lfloor k/d\rfloor$ and $b = k \bmod d$.
Then for every $\sigma \in S_r$ with $p$ cycles and every $h \in S_\alpha$,
\[
\lfloor n/d\rfloor - m_d(\tilde\sigma h) \;\ge\; (r - p)\,a + \lfloor rb/d\rfloor.
\]
\end{corollary}

\begin{proof}
$n = rk = rad + rb$ gives $\lfloor n/d\rfloor = ra + \lfloor rb/d\rfloor$
(the division is exact on $rad$; the leftover $rb$ contributes
$\lfloor rb/d\rfloor$). Subtract the bound of Proposition \ref{prop:bound}.
\end{proof}

\section{Proof of Proposition \ref{prop:main} and the theorem}

\begin{proof}[Proof of Proposition \ref{prop:main}]
Assume $b := k \bmod d \in \{2, \ldots, d-1\}$. We show that
$\lfloor n/d\rfloor - m_d(\tilde\sigma h) \ge 1$ for every $\sigma \in S_r$
and $h \in S_\alpha$, whence every summand of the Molien formula of
Lemma \ref{lem:molien} vanishes at $\zeta_d$ by Lemma \ref{lem:order}.

By Corollary \ref{cor:gap}, the gap is at least
$(r - p)\,a + \lfloor rb/d\rfloor$ where $p = p(\sigma)$ and $a = \lfloor k/d\rfloor$.

\smallskip
\textbf{Case 1: $\sigma = e$.} Then $p = r$, and the gap is at least
$\lfloor rb/d\rfloor$. Since $b \ge 2$, we have $rb \ge 2r = d + 1 > d$,
hence $\lfloor rb/d\rfloor \ge 1$.

\smallskip
\textbf{Case 2: $\sigma \ne e$.} Then $\sigma$ has at least one cycle of
length $\ge 2$, so $p \le r - 1$; hence $r - p \ge 1$. Since $a \ge 0$
and $\lfloor rb/d\rfloor \ge 1$ (as in Case 1),
\[
(r - p)\,a + \lfloor rb/d\rfloor \;\ge\; 0 + 1 \;=\; 1.
\]

In either case, the gap is $\ge 1$, so $P_\sigma(\zeta_d) = 0$.
\end{proof}

\begin{proof}[Proof of the theorem]
Combine Proposition \ref{prop:main} with Lemma \ref{lem:reduce} and the
equivalence $[d]_t \mid f(t) \iff f(\zeta_d) = 0$ (valid since $d$ prime,
$[d]_t = \Phi_d(t)$).
\end{proof}

\section{Discussion}

\begin{remark}[Order of vanishing]
Corollary \ref{cor:gap} in fact shows that every summand of Lemma
\ref{lem:molien} vanishes at $\zeta_d$ to order at least
$(r - p)\,a + \lfloor rb/d\rfloor$. For $\sigma = e$ this is $\lfloor rb/d\rfloor$,
so $P_e(t) = \binom{n}{k^r}_t$ vanishes at $\zeta_d$ to order at least
$\lfloor rb/d\rfloor$. This is in fact the exact order of vanishing
of $\binom{n}{k^r}_t$ at $\zeta_d$ (a direct calculation from
$\binom{n}{k^r}_t = [n]_t!/([k]_t!)^r$: the order is
$\lfloor n/d\rfloor - r\lfloor k/d\rfloor = \lfloor rb/d\rfloor$).
So Corollary \ref{cor:gap} is sharp at $\sigma = e$; term-by-term
vanishing is not enough to detect any \emph{cancellation}, which for
$\sigma = e$ is empty by definition. For $\sigma \ne e$ the bound may
be conservative.
\end{remark}

\begin{remark}[Why the first-period theorem restricted to $k < d$]
The first-period argument of \cite{Clio-first-period} used the stronger
inequality $m_d(\tilde\sigma h) = 0$ (its Cor.\ 4.5), obtained from
$k \le d - 1$: then no cycle of any companion has length divisible by
$d$, so trivially $m_d = 0$. For $k \ge d$ this is false: taking
$\sigma = e$ and $h_i = (12\cdots d)(d+1)\cdots(k)$ gives
$m_d(h) = r \ge 1$. But we do not need $m_d = 0$; the weaker
$m_d < \lfloor n/d\rfloor$ suffices, and Proposition \ref{prop:bound}
shows the weaker inequality throughout the allowed range.
\end{remark}

\begin{remark}[Complementary sharpness at forbidden residues]
When $b \in \{0, 1\}$ (the ``forbidden'' residues), the bound of
Proposition \ref{prop:bound} is achieved with equality at
$\sigma = e$: taking each $h_i \in S_k$ with cycle type
$d^a \cdot 1^b$ gives $m_d(h) = ra$; and
$\lfloor n/d\rfloor - ra = \lfloor rb/d\rfloor$. For $b = 0$: gap $= 0$,
so the $\sigma = e$ Molien term does \emph{not} vanish and there is no
divisibility (matching the empirical pattern). For $b = 1$: $\lfloor r/d\rfloor
= 0$ (since $r < d$), so gap $= 0$ again. Thus the same argument gives
a clean explanation for the ``forbidden'' residues without invoking
cancellation.
\end{remark}

\begin{remark}[Composite $d$]
The prime hypothesis on $d$ is used at two points: (i) to reduce
$[d]_t \mid f$ to $f(\zeta_d) = 0$; (ii) to force $d \mid \ell m
\Rightarrow d \mid m$ for $\ell \le r < d$ prime.
For composite $d = 2r - 1$ (e.g. $r = 5, d = 9$), point (i) becomes
divisibility by each cyclotomic factor $\Phi_e(t)$, $e \mid d$, separately;
point (ii) requires $e > r$ for each prime factor $e$, which fails when
$r$ has a small prime factor $< d$. Thus the argument as stated does
not apply.
\end{remark}

\begin{remark}[Isotypic vs.\ Hilbert-series interpretation]
Setting $\sigma = e$ gives the well-known cyclotomic divisibility
of the $q$-multinomial:
$[d]_t \mid \binom{rk}{k, k, \ldots, k}_t$ iff $b \ge 2$
(equivalently, $b \ne 0, 1$). This is a classical fact and can be
proved directly from
$\binom{rk}{k^r}_t = [rk]_t!/([k]_t!)^r$ by counting cyclotomic factors
of $[m]_t!$ at $\zeta_d$. The content of the present theorem is the
\emph{isotypic} refinement: the same factor $[d]_t$ divides every
$m_\pi(t)$, not merely their $f^\pi$-weighted sum
$\binom{rk}{k^r}_t = \sum_\pi f^\pi m_\pi(t)$.
\end{remark}

\section{Verification}
\label{sec:verify}

The Molien bound of Proposition \ref{prop:bound} was verified by brute
enumeration of $S_\alpha = S_k^r$ for
$(r, k) \in \{(2,5), (3,2)\}$, computing
$\max_h m_d(\tilde\sigma h)$ for each conjugacy class of $\sigma \in S_r$
and checking against the theoretical bound $p \lfloor k/d\rfloor$. In
every case the bound matches the enumerated maximum.

The multinomial vanishing $\binom{n}{k^r}(\zeta_d) = 0$ (i.e., the
$\sigma = e$ case) was verified symbolically for
$(r, k) \in \{(2, 5), (3, 7), (3, 8), (3, 9), (3, 12), (4, 9)\}$.

The full isotypic divisibility $[d]_t \mid m_\pi(t)$ was verified by
explicit computation of $m_\pi(t)$ (via Theorem D expansion plus fake
degrees) for $(r, k, d) = (3, 7, 5)$ and each $\pi \in \{(3), (2,1), (1,1,1)\}$;
in every case $[5]_t \mid m_\pi(t)$. Scripts:
\texttt{\~{}/projects/probes/2026-08-04-prove-second-period/verify.py}.

\begin{thebibliography}{9}
\bibitem{Clio-first-period} C.~Muse (Clio), \emph{$[2r-1]_t$-divisibility
of $S_r$-isotypic multiplicities of $R_{(k^r)}$ for $2 \le k \le 2r-2$:
a Molien-Springer proof}, note (2026-07-31).
\bibitem{Clio-Wreath-general-r} C.~Muse (Clio), \emph{The $r$-wreath
extension of Theorem D and a cyclotomic Hilbert-series identity},
note (2026-07-30).
\bibitem{Macdonald} I.G.~Macdonald, \emph{Symmetric Functions and Hall
Polynomials}, 2nd ed., Oxford Univ.~Press, 1995.
\end{thebibliography}

\end{document}
